\documentclass[12pt]{article}
\usepackage{amsfonts,amsthm,amsmath,amssymb,mathtools}
\usepackage{mathrsfs}
\usepackage{enumitem}
\usepackage{tikz-cd}
\usepackage{geometry}
\usepackage{mlmodern}

\geometry{letterpaper, portrait, margin=1in}

\setcounter{section}{0}

\renewcommand{\thesection}{\S\arabic{section}}
\renewcommand{\thesubsection}{\arabic{section}}
\newcommand{\dd}{\mathop{}\!\mathrm{d}}
\newcommand{\R}{\mathbb{R}}
\newcommand{\Z}{\mathbb{Z}}
\newcommand{\C}{\mathbb{C}}
\newcommand{\Q}{\mathbb{Q}}

\DeclareMathOperator{\lcm}{lcm}
\DeclareMathOperator{\im}{im}

\newtheorem{thm}{Theorem}
\newenvironment{theorem}[1]{
	\renewcommand{\thethm}{#1}
	\begin{thm}
		}{\end{thm}}

\newtheorem{lma}{Lemma}
\newenvironment{lemma}[1]{
	\renewcommand{\thelma}{#1}
	\begin{lma}
		}{\end{lma}}

\theoremstyle{definition}
\newtheorem{ex}{}
\newenvironment{exercise}[1]{
	\renewcommand{\theex}{#1}
	\begin{ex}
		}{\end{ex}}

\title{Project 1: Bremner and MacLeod}
\author{Connor Mooneyhan}
\date{December 9, 2025}

\begin{document}

\maketitle

Prompt: Near the beginning of the class, I pointed out that the equation \begin{equation*}
	\frac{a}{b+c}+ \frac{b}{a+c} + \frac{c}{a+b} = 4
\end{equation*} has positive integer solutions, but the smallest solution is very large (with a, b, c having 81, 79, and 80 digits, respectively).
Another surprising result is that it is not possible to find a positive integer solution to \begin{equation*}
	\frac{a}{b+c}+ \frac{b}{a+c} + \frac{c}{a+b} = N
\end{equation*} if N is odd and positive.
The goal of this project is to prove this and learn any new mathematics needed along the way.
(This result was first proven by Andrew Bremner and Allan MacLeod in the paper ``An unusual cubic representation problem.'')
Their proof relies on properites of quadratic residues, the Jacobi symbol, and the law of quadratic reciprocity.
Learn about these tools, and then give your own accounting of Bremner and MacLeod’s proof.
(You should have the background you need, other than learning about quadratic residues, at the end of Week 11.)

\section{Introduction}
One of the most powerful uses for elliptic curves is in their ability to help us solve equations that often require rational or integer solutions.
This is often done by taking the original equations and mapping them to an appropriate elliptic curve, then studying the rational points on that curve via the Mordell-Weil group, and finally mapping those points back to the original equation.
The following equation is one such example, and it is what we will be studying in this project: \begin{equation}\label{1}
	\frac{a}{b+c}+ \frac{b}{a+c} + \frac{c}{a+b} = N.
\end{equation}

\section{Background}
An integer $q$ is a \textit{quadratic residue mod $p$} if there exists a positive integer $x < p$ such that $x^2 \equiv q \bmod p$.
This allows us to define the Legendre symbol $(n/p)$, where $n$ is an integer and $p$ is a prime, to be 0 if $p$ divides $n$, 1 if $n$ is a quadratic residue mod $p$, and $-1$ otherwise.
The Jacobi symbol is just a generalization of this, and thus we use the same notation.
We define the Jacobi symbol $(n/m)$ for positive odd $m$ to be \begin{equation*}
	\left( \frac{n}{m } \right) = \left( \frac{n}{p_1} \right)^{a_1}\left( \frac{n}{p_2} \right)^{a_2}\cdots\left( \frac{n}{p_k} \right)^{a_k}
\end{equation*}
where $m=p_1^{a_1}p_2^{a_2}\cdots p_k^{a_k}$ is the prime factorization of $m$.
We can now formulate the law of quadratic reciprocity in terms of the Jacobi symbol: \begin{equation*}
	\left( \frac{m }{n} \right)\left( \frac{n}{m } \right) = (-1)^{\frac{(m-1)(n-1)}{4}}
\end{equation*}
for $m$ and $n$ relatively prime odd integers with $n \geq 3$.

\section{Bremner and MacLeod's Proof}
The first major step in the proof is to homogenize \eqref{1}, which includes cancelling out denominators.
This yields a curve for which a cubic model is found of the form \begin{equation*}
	E_N : y^2 = x^3 + (4N^2 + 12N - 3) x^2 + 32(N+3)x.
\end{equation*}
This model comes with maps that will allow translation between points on the curve and solutions to the original equation, and thus the curve $E_N$ will be the primary object of study.
We will focus on positive $N$ for this project.

The first thing we can say about $E_N$ is that if $N > 1$, the discriminant is positive, so it has two components: an unbounded component with $x \geq 0$ and an egg with $x < 0$.
In fact, it is shown that for there to be positive solutions in $a$, $b$, and $c$, the inequality \begin{equation*}
	\frac{1}{2}(3-12N-4N^2-(2N+5)\sqrt{4N^2+4N-15}) < x < -2(N+3)(N+\sqrt{N^2-4})
\end{equation*}
must be satisfied.
This in particular means that all rational points on $E_N$ corresponding to positive solutions of \eqref{1} lie on the egg.
This fact gives rise to what is for our purposes the key theorem: when $N$ is odd, there are no rational points on $E_N$ with $x < 0$.

To prove this, we assume $N$ to be odd and let $M = (N+3)/2$.
Substituting this into $E_N$, we get \begin{equation*}
	E_M : y^2 = x(x^2 + (16M^2 - 24M-3)x + 64M).
\end{equation*}
A point on this curve is found to have the form $x=dr^2/s^2$ where $d, r, s \in \mathbb{Z}$, $d$ is squarefree, $r$ and $s$ share no common factors, and $d$ divides $64M$.
This gives us the quartic \begin{equation}\label{2}
	dr^4 + (16M^2 - 24M - 3)r^2s^2 + \frac{64M}{d}s^4.
\end{equation}
The rest of the proof then reduces to showing that \eqref{2} has no solutions $r, s$ where $d < 0$.
This is treated in two cases: when $d$ is odd and when $d$ is even.

\section{The First Case}
We first consider when $d < 0$ is odd.
Let $u > 0$ be odd as well and $M = um$.
Then \eqref{2} becomes \begin{equation*}
	-ur^4 + (16M^2 - 24M - 3)r^2s^2 - 64ms^4.
\end{equation*}
Here we calculate a Jacobi symbol: \begin{align*}
	\left( \frac{-u}{4M-1} \right) & = \left( \frac{-1}{4M-1} \right) \left( \frac{u}{4M-1} \right) \\
	                               & = -\left( \frac{4M-1}{u} \right) (-1)^{(u-1)/2}                \\
	                               & = -\left( \frac{-1}{u} \right)(-1)^{(u-1)/2}                   \\
	                               & = -1.
\end{align*}
We use this to deduce a contradiction.
Indeed, if $\left( \frac{-u}{p} \right) = 1$ for every prime $p$ dividing $4M - 1$, then we would have $\left( \frac{-u}{4M-1} \right) = 1$ because it would be a product $\left( \frac{-u}{p_1} \right)\cdots \left( \frac{-u}{p_n} \right)$ where each $p_i$ is a prime dividing $4M-1$.
This contradicts the above calculation.
Thus there must exist a prime $p$ dividing $4M-1$ for which $\left( \frac{-u}{p} \right) = -1$.
Then we have \begin{align*}
	-2ur^2 + (16M^2 - 24M - 3)s^2 & \equiv 0 \bmod p      \\
	-2ur^2 - 8s^2                 & \equiv 0 \bmod p      \\
	4s^2                          & \equiv -ur^2 \bmod p.
\end{align*}
Then we must have $r \equiv s \equiv 0 \bmod p$, which is a contradiction.

\section{The Second Case}
Now, if $d < 0$ is even, then we can let $d = -2u$ where $u > 0$ is odd and again let $M = um$.
Then \eqref{2} becomes \begin{equation*}
	-2ur^4 + (16M^2 - 24M - 3)r^2s^2 - 32ms^4.
\end{equation*}
This now breaks down into a couple of subcases, namely when $M$ is even and when $M$ is odd.

If $M$ is even, then \begin{align*}
	\left( \frac{-2u}{4M-1} \right) & = \left( \frac{-2}{4M-1} \right)\left( \frac{4M-1}{u} \right)(-1)^{(u-1)/2} \\
	                                & = -\left( \frac{-1}{u} \right)(-1)^{(u-1)/2} = -1.
\end{align*}
Similar to before, we can deduce that there exists a prime $p$ dividing $4M - 1$ where $\left( \frac{-2u}{p} \right) = -1$.
Then we get a similar implication chain to before, where \begin{align*}
	-4ur^2 + (16M^2 - 24M - 3)s^2 & \equiv 0 \bmod p       \\
	-4ur^2 - 8s^2                 & \equiv 0 \bmod p       \\
	4s^2                          & \equiv -2ur^2 \bmod p.
\end{align*}
Again, this leads to the contradiction $r \equiv s \equiv 0 \bmod p$.

Now if $M$ is odd, then $m$ is also odd.
We reduce \eqref{2} to the following expression (which has been assigned to a variable $X$): \begin{equation*}
	-2ur^4 + (16M^2 - 24M - 3)r^2s^2 - 32ms^4 = X.
\end{equation*}
This would imply that $s$ is odd.
As for $r$, it cannot be odd modulo 4, so it is even.
Then we have $-3(r/2)^2 \equiv X \bmod 8$, implying that $r/2$ is even, and then \begin{align*}
	-3(r/4)^2 - 2m & \equiv X \bmod 4 \\
	-3(r/4)^2 - 2  & \equiv X \bmod 4,
\end{align*}
which is impossible.

\section{Conclusion}
We have now shown that for odd $N$, any rational points on the elliptic curve $E_N$ corresponding to positive rational solutions to \eqref{1} must be on the egg, but in fact there are no points on the egg, so there are no such solutions.

\section{Bibliography}
https://mathworld.wolfram.com/QuadraticResidue.html
\\
Used for learning about quadratic residues.
\\\\
https://mathworld.wolfram.com/JacobiSymbol.html
\\
Used for learning about the Jacobi symbol and the law of quadratic reciprocity.
\\\\
https://mathworld.wolfram.com/LegendreSymbol.html
\\
Used for learning about the Legendre symbol.

\end{document}
